\chapter{Expressiveness, Translation, and Representation Tables}
\label{sec:p14-expressiveness-and-representation}
\index{translation!reference}\index{representation!reference}

A translation maps specified syntax to specified syntax; a representation
maps mathematical objects and, when asserted, their morphisms. Its guarantee
may concern truth at a point, consequence, theoremhood, satisfiability,
equations or reconstruction. These are distinct properties. The tables below
state the objects being mapped, the forward and backward conclusions, and
the assumptions needed to obtain them. Their proof links retain the actual
language, calculus, source presentation and model regime.

The general expressiveness map in \cref{sec:p14-expressive-table} locates
systems. The present chapter gives their mapping interfaces and conditions.
For the decision problem and encoding behind a complexity statement use
\cref{sec:p14-decision-table}; an expressive embedding by itself supplies
no size, time or proof-length bound \cite{Enderton2001,Libkin2004,
BlackburnDeRijkeVenema2001}.

\section{Embeddings and faithful translations between systems}
\label{sec:p14-translation-embeddings}
\subsection{Expressive-power map}
\label{sec:p14-expressive-table}

\begin{longtable}{@{}p{0.21\textwidth}p{0.27\textwidth}p{0.42\textwidth}@{}}
\caption{Expressive-power map}\\
\toprule
Comparison & Inclusion or equivalence & Witnessing boundary \\
\midrule
\endfirsthead
\caption[]{Expressive-power map (continued)}\\
\toprule
Comparison & Inclusion or equivalence & Witnessing boundary \\
\midrule
\endhead
\midrule\multicolumn{3}{r}{\emph{continued on next page}}\\
\endfoot
\endlastfoot
modal logic vs first-order logic & standard translation embeds modal logic in FO & bisimulation-invariant boundary \cite{VanBenthem1983} \\
CPL vs first-order logic & propositional formulas embed as nullary-predicate cases & domain-sensitive quantification \cite{Enderton2001} \\
FO vs existential SO on finite structures & proper inclusion & ESO captures NP \cite{Fagin1974} \\
FO vs least fixed-point logic on ordered finite structures & proper inclusion & LFP captures P \cite{Immerman1986} \\
LFP vs ESO on ordered finite structures & unconditional inclusion & \(\mathsf P\subseteq\mathsf{NP}\); equality exactly when \(\mathsf P=\mathsf{NP}\) \cite{Immerman1999} \\
LTL vs FO on words & LTL with past captures FO point properties & Kamp theorem \cite{Kamp1968} \\
CTL vs LTL & incomparable & branching/path witnesses \cite{EmersonHalpern1986} \\
CTL and LTL vs CTL* & both proper fragments & unrestricted state/path nesting \cite{EmersonHalpern1986} \\
modal \(\mu\)-calculus vs modal logic & proper extension & fixed-point reachability \cite{Kozen1983} \\
full SO vs Henkin SO & same syntax, different validities & full versus general relation ranges \cite{Henkin1950,Shapiro1991} \\
\bottomrule
\end{longtable}

All comparisons hold over the displayed structure class and semantics. A
translation that preserves satisfiability need not preserve formula size,
validity, proof length, or intended interpretation.


\subsection{A satisfaction interface for the modal standard translation}
\label{sec:p14-translation-standard-interface}

The basic modal language uses propositional atoms, Boolean connectives and
\(\Box\), with \(\Diamond A=\neg\Box\neg A\). The first-order target
has a binary predicate \(R\), a distinct unary predicate \(P_p\) for
every modal atom \(p\), and a displayed free world variable \(x\).
The exact map \(\operatorname{ST}_x\) is
\cref{def:p08-modal-standard-translation}; its characteristic clauses are
\[
 \operatorname{ST}_x(p)=P_p(x),\qquad
 \operatorname{ST}_x(\Box A)=
       \forall y(Rxy\to\operatorname{ST}_y(A)),
\]
\[
 \operatorname{ST}_x(\Diamond A)=
       \exists y(Rxy\land\operatorname{ST}_y(A)),
\]
using fresh bound variables. It commutes with Boolean operations and
preserves falsity. This is a formula transformation with a free world
parameter, not an identification of modal syntax with every first-order
formula about a relation \cite{BlackburnDeRijkeVenema2001}.

\begin{hbdefinition}{def:p14-standard-translation-models}{The two model transformations}
For a nonempty Kripke structure \(\mathcal M=(W,R,V)\), let
\(\mathfrak A_{\mathcal M}\) be the first-order structure with domain
\(W\), its same binary relation, and
\(P_p^{\mathfrak A_{\mathcal M}}=V(p)\). Conversely, a nonempty
first-order structure \(\mathfrak A\) in this signature gives the Kripke
structure \(\mathcal M_{\mathfrak A}\) with that domain and relation,
and \(V(p)=P_p^{\mathfrak A}\).

For a pointed model \((\mathcal M,w)\), use an object-variable assignment
\(s\) with \(s(x)=w\). Other variable values do not affect a translated
formula's truth because its only possible free variable is \(x\).
For a set \(\Gamma\) of modal formulas,
\(\operatorname{ST}_x[\Gamma]\) is the memberwise image. Local modal
consequence requires all premises and the conclusion at the same world.
\end{hbdefinition}

\begin{hbtheorem}{thm:p14-standard-translation-interface}{proved=annex:p14.translation.standard-interface}{Truth, local consequence, satisfiability and validity}
For every basic modal formula \(A\), pointed nonempty model
\((\mathcal M,w)\), and assignment \(s(x)=w\),
\[
 \mathcal M,w\entails A
       \quad\Longleftrightarrow\quad
 \mathfrak A_{\mathcal M},s\entails\operatorname{ST}_x(A).
\]
For every formula set \(\Gamma\), local consequence over all nonempty
Kripke frames is equivalent to ordinary first-order consequence
\[
 \Gamma\entails^{\ell}_{\mathrm{Kripke}} A
       \quad\Longleftrightarrow\quad
 \operatorname{ST}_x[\Gamma]\entails\operatorname{ST}_x(A).
\]
Moreover, modal satisfiability of \(A\) is equivalent to first-order
satisfiability of \(\exists x\operatorname{ST}_x(A)\), and modal validity
is equivalent to first-order validity of
\(\forall x\operatorname{ST}_x(A)\). These statements concern the given
model transformations and the declared formula language. The world domain
need not be finite or countable, and no Choice assumption is used.
They do not assert that every first-order formula is a modal translation.
\end{hbtheorem}

On the Boolean fragment this gives a propositional embedding by
\(p\mapsto P_p(x)\). A singleton first-order domain realizes any
propositional valuation; evaluating at a world recovers its valuation.
Where the signature permits nullary predicates, the separate direct map
\(p\mapsto P_p\) gives a closed-formula presentation instead. Neither
representation adds domain-sensitive quantification to the source language.

\subsection{Classical, intuitionistic and modal proof translations}

In the following negative translations \(P\) denotes a genuine atomic
formula and excludes the grammar constant \(\bot\). Negation in the
target is \(A\to\bot\). The recursive maps act memberwise on a context
when that is explicitly stated. The clauses and result standing are fixed
by \cref{def:p08-negative-translation-maps,tab:p08-negative-translations}
and the inherited G\"odel--Gentzen entry.

\begin{longtable}{@{}p{.19\textwidth}p{.35\textwidth}p{.36\textwidth}@{}}
\caption{Classical, intuitionistic and modal proof translations}\\
\toprule
Source and target & Actual transformation & Forward and backward guarantee, with conditions \\
\midrule
\endfirsthead
\caption[]{Classical, intuitionistic and modal proof translations (continued)}\\
\toprule
Source and target & Actual transformation & Forward and backward guarantee, with conditions \\
\midrule
\endhead
\midrule\multicolumn{3}{r}{\emph{continued on next page}}\\
\endfoot
\endlastfoot
Propositional \(LK\) to \(LJ\), Glivenko &
For the input query \((\Gamma,A)\), leave \(\Gamma\) unchanged and
replace the conclusion by \(\neg\neg A\). &
\(LK:\Gamma\Rightarrow A\) iff \(LJ:\Gamma\Rightarrow\neg\neg A\),
for finite multisets and propositional rules only. This is the exact
context form of \cref{thm:glivenko}; it asserts no quantifier extension. \\
First-order \(LK\) to \(LJ\), G\"odel--Gentzen &
\(P^N=\neg\neg P\), \(\bot^N=\bot\); commute with
\(\land,\to,\forall\); use
\((A\lor B)^N=\neg(\neg A^N\land\neg B^N)\) and
\((\exists xA)^N=\neg\forall x\neg A^N\). &
\(LK:\Gamma\Rightarrow A\) iff \(LJ:\Gamma^N\Rightarrow A^N\).
Translate every premise; the falsity clause has priority over the atom
clause. The fixed inherited presentation is \cref{thm:negative-translation}. \\
First-order \(LK\) to \(LJ\), Kolmogorov &
Double-negate atoms, falsity, and every result of a connective or
quantifier construction after recursively translating its arguments. &
\(\Gamma\derives_{LK}A\) iff
\(\Gamma^{\mathrm{Ko}}\derives_{LJ}A^{\mathrm{Ko}}\),
for finite contexts. The exact \(\bot\), Boolean and quantifier clauses
are \cref{def:p08-negative-translation-maps}; the result is
\cref{thm:p08-kolmogorov-embedding}. \\
First-order \(LK\) to \(LJ\), Kuroda &
The inner map commutes with Boolean operations and \(\exists\), leaves
atoms and falsity unchanged, and maps \(\forall xA\) to
\(\forall x\neg\neg A_{\mathrm{Ku}}\). Then wrap the whole result in
double negation. &
\(\Gamma\derives_{LK}A\) iff
\(\Gamma^{\mathrm{Ku}}\derives_{LJ}A^{\mathrm{Ku}}\), for finite
contexts. The target here is intuitionistic logic, with its falsity rule.
\cref{thm:p08-kuroda-embedding} makes no minimal-logic claim. \\
IPL to \(S4\) &
\(p^{\Box}=\Box p\), \(\bot^{\Box}=\bot\); commute with
\(\land,\lor\); map implication to
\(\Box(A^{\Box}\to B^{\Box})\). &
\(\derives_{\mathrm{IPL}}A\) iff \(\derives_{S4}A^{\Box}\),
the theoremhood equivalence in \cref{thm:p08-ipl-s4-embedding}.
The cited result does not by itself supply arbitrary-premise consequence
or a size bound. \\
Nonvacuous and vacuous \(NJ\) &
Keep formulas and open assumptions. To populate an empty discharge class,
the nonvacuous target inserts conjunction introduction followed by projection. &
Both directions preserve conclusions and open assumption occurrences,
\cref{thm:p07-nj-vacuous-consequence-bridge}. The transformation does not
preserve normality; \cref{thm:p07-nj-nonvacuous-obstruction} gives the
reduction boundary. \\
\bottomrule
\end{longtable}

The four negative maps have separately specified placement of negations.
For example the G\"odel--Gentzen falsity image is \(\bot\), whereas
Kolmogorov and the full Kuroda map give \(\neg\neg\bot\).
These expressions can be equivalent in a target system without being the
same recursion. The exact G\"odel--Gentzen statement uses the source's
defined-negation reading; a native primitive-negation presentation needs
its own correspondence. The listed results retain their attributed
presentations \cite{TroelstraVanDalen1988,Glivenko1929,
Kolmogorov1925ExcludedMiddle,Kuroda1951,Godel1933}.

\subsection{Sorts, incidence and fresh symbols}

\begin{longtable}{@{}p{.19\textwidth}p{.35\textwidth}p{.36\textwidth}@{}}
\caption{Sorts, incidence and fresh symbols}\\
\toprule
Source and target & Mapping of syntax and models & Guarantee and retained conditions \\
\midrule
\endfirsthead
\caption[]{Sorts, incidence and fresh symbols (continued)}\\
\toprule
Source and target & Mapping of syntax and models & Guarantee and retained conditions \\
\midrule
\endhead
\midrule\multicolumn{3}{r}{\emph{continued on next page}}\\
\endfoot
\endlastfoot
Strict many-sorted to one-sorted FO &
Mark each sort by \(Q_{\mathsf s}\); replace
\(\forall x{:}\mathsf s\,A\) by
\(\forall x(Q_{\mathsf s}(x)\to\operatorname{tr}_{\mathrm{MS}}(A))\),
and existential quantification by the corresponding guarded conjunction. &
The consequence equivalence includes the nonempty/disjoint sort and symbol
typing axioms and the total sorted-assignment guard.
\cref{def:p02-many-sorted-translation,def:p02-ms-variable-guard,
thm:p02-many-sorted-translation}. Arbitrary unguarded one-sorted assignments
are not the matching sorted assignments. \\
Relational Henkin SOL to many-sorted FO &
Object sort \(D\), relation sort \(S_n\), and incidence
\(X^n(\bar t)\mapsto E_n(X^n,\bar t)\). Predicate constants become
relation-sort constants \(\widehat P\); relation quantifiers become
quantifiers of their relation sort. Tag all domains and assignments. &
Satisfaction is equivalent for the paired structures and assignments,
\cref{thm:p02-sol-incidence-satisfaction}. The exact eligible
comprehension/choice translations characterize the matched model class.
Under ambient AC, the translation-defined calculus is strongly adequate,
\cref{thm:p02-sol-derived-adequacy}. \\
Closed FO formulas to clauses &
Simplify, take NNF, miniscope, rectify, introduce fresh witness functions,
then form the universal prefix and distributive clause matrix. &
The Boolean and rectification stages preserve equivalence. Skolem steps
give expansion/reduct satisfiability, with Choice for arbitrary-domain
expansion and freshness outside the declared signature.
\cref{thm:p12-small-rectification,thm:p12-small-skolem,
thm:p12-small-cnf}. Same-signature equivalence is not the guarantee. \\
\bottomrule
\end{longtable}

The incidence translation changes the application predicate and quantifier
sorts. Erasing the relation binder would not perform that transformation.
The source's cited typed completeness theorem in
\cref{sec:simple-type-theory-02}
belongs to Henkin's specified typed system and general models
\cite{Henkin1950}. It is distinct from the relational translation-defined
calculus. Admissibility of the displayed relational rules is not a proof
that those rules generate every derivation of the translated system.
The exact boundary is \cref{def:p02-sol-native-boundary}.

Clausification's model relation likewise matters. If a formula only asserts
\(\exists xP(x)\), naming a fresh witness constant chooses an expansion;
it does not make \(P(c)\) true for every preassigned value of a constant
already in the signature. The direct and optimized constructions retain
their separate source presentations
\cite{NonnengartWeidenbach2001SmallClauses}.

\section{Frame correspondences and algebraic semantics}
\label{sec:p14-translation-frame-algebra}
\subsection{Modal axiom--frame--algebra correspondences}
\label{sec:p14-modal-correspondence}

The rows use the normal-modal correspondence and complex-algebra translations
in \cite{BlackburnDeRijkeVenema2001,Goldblatt1989}.

\begin{tabularx}{\textwidth}{@{}l l l X@{}}
\toprule
Axiom & Frame condition & Algebraic inequality & System use \\
\midrule
\(\Box(p\to q)\to(\Box p\to\Box q)\) & none beyond normality & \(\Box\) preserves finite meets and \(1\) & \(\mathsf K\) \\
\(\Box p\to p\) & reflexive & \(\Box a\le a\) & \(\mathsf T\) \\
\(\Box p\to\Box\Box p\) & transitive & \(\Box a\le\Box\Box a\) & axiom 4 \\
\(p\to\Box\Diamond p\) & symmetric & \(a\le\Box\Diamond a\) & axiom B \\
\(\Box p\to\Diamond p\) & serial & \(\Box0=0\) & axiom D \\
\(\Diamond p\to\Box\Diamond p\) & Euclidean & \(\Diamond a\le\Box\Diamond a\) & axiom 5 \\
T + 4 & reflexive, transitive & interior algebra & \(\mathsf{S4}\) \\
T + 5 & equivalence relation & monadic Boolean algebra & \(\mathsf{S5}\) \\
\bottomrule
\end{tabularx}

The table records standard correspondence for normal propositional modal
logic. It does not transfer unchanged to neighborhood, many-valued, first-order
modal, dynamic, or nonnormal semantics.


For a normal propositional modal frame \((W,R)\), its full complex algebra
has carrier \(\mathcal P(W)\), Boolean set operations and
\[
 \Box U=\{w:\forall v(wRv\Longrightarrow v\in U)\},\qquad
 \Diamond U=\{w:\exists v(wRv\land v\in U)\}.
\]
Evaluating a formula in this algebra returns exactly its set of satisfying
worlds. Validity at every world under every propositional valuation becomes
equality to the top element under every algebra valuation. The constructions
are \cref{def:p10-bao-complex-algebra,def:p10-bao}; the frame-schema
correspondences are \cref{thm:p08-modal-correspondence}
\cite{Chellas1980,Venema2005Preprint}.

\begin{longtable}{@{}p{.27\textwidth}p{.28\textwidth}p{.35\textwidth}@{}}
\caption{Frame correspondences and algebraic semantics}\\
\toprule
Schema valid under every valuation & Exact frame condition & Algebraic form in the full complex algebra \\
\midrule
\endfirsthead
\caption[]{Frame correspondences and algebraic semantics (continued)}\\
\toprule
Schema valid under every valuation & Exact frame condition & Algebraic form in the full complex algebra \\
\midrule
\endhead
\midrule\multicolumn{3}{r}{\emph{continued on next page}}\\
\endfoot
\endlastfoot
\(\Box(p\to q)\to(\Box p\to\Box q)\) &
Normal Kripke semantics; no further restriction on \(R\). &
\(\Box\) preserves finite meets and \(1\). \\
\(\Box p\to p\), T & Reflexive \(R\). & \(\Box a\leq a\). \\
\(\Box p\to\Diamond p\), D & Serial \(R\). & \(\Box0=0\). \\
\(\Box p\to\Box\Box p\), 4 & Transitive \(R\). &
\(\Box a\leq\Box\Box a\). \\
\(p\to\Box\Diamond p\), B & Symmetric \(R\). &
\(a\leq\Box\Diamond a\). \\
\(\Diamond p\to\Box\Diamond p\), 5 & Euclidean \(R\). &
\(\Diamond a\leq\Box\Diamond a\). \\
\bottomrule
\end{longtable}

The element variable \(a\) ranges over every algebra element. Testing
one formula valuation on a frame does not establish its frame condition.
The complete separating frames in \cref{thm:p08-mce-separations}
show that seriality, symmetry and transitivity cannot be substituted for
reflexivity. The B schema alone gives symmetry; the named \(KTB\) system
also requires T. The table's equivalences concern ordinary Kripke frames
and their full powerset algebras, with the source's normal semantics.

\paragraph{Sahlqvist correspondence and its syntactic boundary.}
Definition~\ref{def:p08-sahlqvist-class} fixes the displayed class:
antecedents are built from boxed atoms, negative formulas and the Boolean
constants using conjunction, disjunction and diamond; consequents are
positive, and finite conjunctions of the resulting implications are allowed.
For that class, Theorem~\ref{thm:p08-sahlqvist} cites effective first-order
correspondence, canonicity and completeness for the corresponding elementary
frame class to \cite{Sahlqvist1975}.  A schema outside this sufficient class
is not thereby non-correspondent: the distinct source account and its
non-Sahlqvist example are
Theorems~\ref{thm:p08-bvb501-sahlqvist-correspondence}
and~\ref{thm:p08-bvb501-first-order-boundaries}.  Their original proofs
remain source-cited, not newly proved by this lookup paragraph.

\subsection{Object and arrow representations}

\begin{longtable}{@{}p{.20\textwidth}p{.35\textwidth}p{.35\textwidth}@{}}
\caption{Object and arrow representations}\\
\toprule
Structures being compared & Explicit map & Forward/backward property and its domain \\
\midrule
\endfirsthead
\caption[]{Object and arrow representations (continued)}\\
\toprule
Structures being compared & Explicit map & Forward/backward property and its domain \\
\midrule
\endhead
\midrule\multicolumn{3}{r}{\emph{continued on next page}}\\
\endfoot
\endlastfoot
Boolean algebra and Stone clopens &
\(B\to\operatorname{Clop}(X_B)\), \(a\mapsto\widehat a\), where
\(X_B\) is the ultrafilter space and \(\widehat a=\{U:a\in U\}\). &
Under BPIT this is a Boolean isomorphism onto all clopens,
\cref{thm:p10-stone-representation}. A Stone space is recovered by
\(x\mapsto\{C:x\in C\}\), \cref{thm:p10-stone-space-reconstruction}.
Trivial algebras and the empty Stone space are included. \\
Stone duality on arrows &
For \(h:B\to C\), use \(X(h):X_C\to X_B\),
\(U\mapsto h^{-1}[U]\). For a continuous \(g:X\to Y\), use
\(g^{-1}:\operatorname{Clop}(Y)\to\operatorname{Clop}(X)\). &
The functors reverse arrows and satisfy the stated natural reconstruction
identities under BPIT, \cref{thm:p10-stone-duality}. An object isomorphism
alone would not supply this arrow correspondence. \\
Boolean algebra with finitary operators and its canonical complex algebra &
\(\beta_B:B\to(\mathbf B^\bullet)^+\),
\(b\mapsto\widehat b\), with the exact ultrafilter relation for each
operation's arity. &
BPIT-qualified embedding preserves every operator, including nullary
operations, \cref{thm:p10-jonsson-tarski}. Its image is the clopen algebra;
it need not be all subsets. Normality and coordinatewise additivity are
the operator hypotheses. \\
IPL formulas and Heyting elements &
Interpret \(\bot,\land,\lor,\to\) by \(0,\wedge,\vee,\Rightarrow\)
and atoms by a valuation. The canonical map sends a formula to its
interderivability class. &
Theoremhood iff value one under every Heyting valuation; for finite
premises, derivability iff the meet of premise values is below the
conclusion. \cref{thm:p10-heyting-completeness}. The empty meet is one. \\
Arrow-only STLC and its context category &
Types denote CCC objects; application denotes evaluation and abstraction
transposition. Finite contexts are products; term tuples modulo
\(\beta\eta\) form arrows composed by substitution. &
Typed \(\beta\eta\) equality iff denotations agree in every CCC
interpretation, \cref{thm:p10-stlc-ccc-completeness}. Context products add
no product or unit term constructor to the registered arrow-only syntax.
\cref{thm:p10-stlc-context-ccc}. \\
\bottomrule
\end{longtable}

The Stone/Boolean-operator source locators distinguish the
Burris--Sankappanavar 2012 text and Venema's 19 October 2005 preprint from
the separately recorded published Venema chapter
\cite{BurrisSankappanavar2012,Venema2005Preprint,Venema2007}.
The STLC interpretation and context construction use the Pitts draft dated
18 May 1995, with the arrow-only application proved in the handbook
\cite{Pitts1995Draft}. These are distinct presentations, rather than
one undated categorical correspondence.

\section{Standard, Henkin, finite, and other model regimes}
\label{sec:p14-translation-model-regimes}

\begin{longtable}{@{}p{.21\textwidth}p{.34\textwidth}p{.35\textwidth}@{}}
\caption{Standard, Henkin, finite, and other model regimes}\\
\toprule
Regime & What ranges over what & Applicable guarantee and boundary \\
\midrule
\endfirsthead
\caption[]{Standard, Henkin, finite, and other model regimes (continued)}\\
\toprule
Regime & What ranges over what & Applicable guarantee and boundary \\
\midrule
\endhead
\midrule\multicolumn{3}{r}{\emph{continued on next page}}\\
\endfoot
\endlastfoot
Unrestricted first-order structures & Object quantifiers range over the
declared nonempty domain; all structures satisfying the background theory
are considered. & Soundness uses these semantics. The general-signature
strong completeness and compactness entries assume ambient AC; countable-
signature completeness has the separate Choice-free proof. No finite-model
restriction is implicit.
\cref{thm:fol-soundness,thm:fol-completeness,
cor:fol-completeness-countable,thm:fol-compactness}. \\
Finite first-order structures & Restrict those object domains to finite
ones while retaining the satisfaction clauses. & Finite compactness fails;
on the stated binary-relation signatures finite validity is not enumerable.
\cref{def:p03-finite-compactness-failure,thm:trakhtenbrot}. This is
different from evaluating one given finite model. \\
Full relational second-order structures & For each arity \(n\geq1\),
relation variables range over every subset of \(M^n\). & Full
comprehension, full \(\mathsf{PA}_2\) categoricity and the full
noncompactness witness use this range.
\cref{thm:p02-sol-full-comprehension,thm:p02-pa2-categoricity,
thm:p02-sol-full-noncompact}. General validity of graph-choice instances
additionally depends on metatheoretic Choice. \\
Arbitrary relational Henkin structures & Each relation domain is a
specified nonempty \(\mathcal R_n\subseteq\mathcal P(M^n)\);
the remaining components interpret nonlogical symbols. & Relation-variable
quantifiers range over
those selected objects. Comprehension or selectors do not follow from
nonemptiness alone. \cref{def:p02-sol-henkin-structure}. \\
Matched relational Henkin structures & Predicate constants belong to the
appropriate selected domains, which satisfy exactly the eligible
comprehension and graph-choice schema families. &
The incidence/guard reduction gives the AC-qualified strong adequacy of
the translation-defined calculus. \cref{def:p02-sol-matched-henkin,
thm:p02-sol-derived-adequacy}. Native rule generation is a further question. \\
Standard simple-type structures &
\(D_o=\{0,1\}\); every arrow type has full function space
\(D_{\alpha\to_{\mathrm{ty}}\beta}=D_\beta^{D_\alpha}\). &
Application and quantification use these full functions. This is
the standard-structure definition in \cref{sec:simple-type-theory-02},
not the selected-domain completeness
claim of Henkin's typed source system. \\
Henkin simple-type structures & Choose nonempty function subsets closed
under denotations of every well-typed lambda abstraction and assignment. &
Closure makes abstraction denotation defined; quantifiers range only over
selected functions. The separate Henkin-structure definition and cited
typed completeness statement are in \cref{sec:simple-type-theory-02}.
Selecting arbitrary nonempty function
sets without closure is insufficient. \\
Modal general frames & Variable valuations range over a declared algebra
of admissible subsets closed under the required operations. & Under BPIT,
every normal logic has the given algebraic and general-frame completeness.
Passing to every subset is an additional canonicity condition.
\cref{thm:p10-modal-general-frame-strong,thm:p10-modal-canonicity-completeness}. \\
Ordered finite descriptive structures & Finite relational structures have
a distinguished total order and the stated finite encoding. &
FO(LFP) captures \(\mathsf P\), and ESO captures \(\mathsf{NP}\).
\cref{thm:p04-immerman-vardi,thm:p03-fagin}. Unordered capture does not
follow by deleting the order from a theorem's hypotheses. \\
\bottomrule
\end{longtable}

The two capture theorems yield containment of FO(LFP)-definable properties
in ESO-definable properties on the ordered finite regime, because every
deterministic polynomial computation is also a nondeterministic polynomial
computation. Equality of those complete property classes is equivalent to
\(\mathsf P=\mathsf{NP}\). The capture statements alone do not establish
strictness or equality. A claim about an individual fixed vocabulary or
restricted fragment also needs that vocabulary or restriction specified.

For temporal logic, the comparison uses another typing convention:
LTL formulas are path formulas, whereas CTL formulas are state formulas.
The state lift \(A\varphi\) universally quantifies the LTL path formula
over paths from the state. On finite total transition structures,
\(\mathsf{AG}\mathsf{EF}p\) is the cited CTL separation from state-lifted
LTL, and \(A(\mathsf F\mathsf Gp)\) is the reverse separation from CTL.
CTL* embeds both with state/path typing retained, as stated in
\cref{thm:p09-ltl-ctl-boundary}\cite{EmersonHalpern1986}.
An untyped assertion that LTL is a fragment of CTL would lose this comparison's
model and formula conditions.

\section{Definability, interpolation, preservation, and invariance}
\label{sec:p14-translation-definability-invariance}

\begin{longtable}{@{}p{.21\textwidth}p{.35\textwidth}p{.34\textwidth}@{}}
\caption{Definability, interpolation, preservation, and invariance}\\
\toprule
Property being compared & Exact forward/backward statement & Model, vocabulary and proof conditions \\
\midrule
\endfirsthead
\caption[]{Definability, interpolation, preservation, and invariance (continued)}\\
\toprule
Property being compared & Exact forward/backward statement & Model, vocabulary and proof conditions \\
\midrule
\endhead
\midrule\multicolumn{3}{r}{\emph{continued on next page}}\\
\endfoot
\endlastfoot
Craig interpolation & From valid \(\varphi\to\psi\), obtain a shared-
vocabulary \(\chi\) with valid \(\varphi\to\chi\) and
\(\chi\to\psi\). & The sentence form is
\cref{thm:p02-craig-interpolation}; constants, functions and predicates
must occur on both sides. Logical equality is exempt. Finite-model
interpolation is a different claim. \\
Beth definability & Uniqueness of the interpreted predicate over the same
reduct language implies an explicit formula in that language; an explicit
definition plainly ensures that uniqueness. &
\cref{thm:p03-beth-definability} supplies the ambient-Choice proof for
possibly infinite first-order theories/signatures and the finite-theory
support conclusion. The original sentence/predicate entry
\cref{thm:p02-beth-definability} retains its citation. \\
Substructure preservation & An elementary class is closed under
substructures iff axiomatizable by universal sentences. & All nonempty
first-order structures, \cref{thm:p03-los-tarski}. The finite restriction
does not satisfy this classical equivalence in general. \\
Extension preservation & A sentence is preserved from a substructure
to an extension iff equivalent to an existential sentence. & The
sentence-level dual in \cref{thm:p03-los-tarski}; keep the signature
and extension relation. Mere embedding is not elementary embedding. \\
Chain-union preservation & An elementary class is closed under the given
substructure-chain unions iff axiomatizable by
\(\forall\bar x\exists\bar y\theta\) with quantifier-free \(\theta\). &
Nonempty ordinal-indexed chains and compatible unions, not only elementary
chains. \cref{thm:p03-chang-los-suszko}. Empty quantifier blocks are allowed. \\
Lyndon positivity & Preservation under surjective homomorphisms iff
equivalence to a positive formula. & Relational signature with equality,
all structures, and positive syntax using \(\land,\lor,\forall,\exists\).
\cref{thm:p03-lyndon-preservation}. Arbitrary homomorphisms and
existential-positive syntax are a distinct pairing. \\
Modal invariance & Bisimilar pointed models agree on every basic modal
formula. Conversely, a first-order formula with one free world variable
is a standard translation iff bisimulation invariant. &
\cref{thm:p08-modal-bisimulation-invariance,thm:p08-van-benthem} concern
all Kripke structures. A finite-model characterization would require its
own theorem; the syntax also keeps the displayed world parameter. \\
Finite-rank first-order invariance & Winning the \(k\)-round EF game
iff agreement on sentences of quantifier rank at most \(k\). & A fixed
finite relational vocabulary, \cref{thm:p03-ef}. Inexpressibility
requires a separating inside/outside pair for every rank,
\cref{thm:p03-ef-inexpressibility}. \\
\bottomrule
\end{longtable}

The cited preservation statements retain their original quantifiers and
map types \cite{ChangKeisler1990,Lyndon1959}. The model-theoretic
Craig--Robinson--Beth route and the separately supplied Maehara construction
have their own proof links; neither permits replacing the common language
by the union language. In \cref{sec:p03-beth-craig}, for instance,
\(A=P\land Q\) and \(B=P\lor R\) share only the nonlogical symbol
\(P\). A definition using \(Q\) in the union language is not thereby
an interpolant. Finite Craig and Beth failures at
\cref{thm:p09-finite-craig-beth} retain their cited exercise-level source
detail and do not acquire a full proof from this comparison.

\subsection{Representing uncertainty without changing its question}

\begin{longtable}{@{}p{.21\textwidth}p{.35\textwidth}p{.34\textwidth}@{}}
\caption{Representing uncertainty without changing its question}\\
\toprule
Objects mapped & Explicit operation & Preserved quantity and limitation \\
\midrule
\endfirsthead
\caption[]{Representing uncertainty without changing its question (continued)}\\
\toprule
Objects mapped & Explicit operation & Preserved quantity and limitation \\
\midrule
\endhead
\midrule\multicolumn{3}{r}{\emph{continued on next page}}\\
\endfoot
\endlastfoot
Probability law and changed coordinates & For a bimeasurable bijection
\(S\), push the reference measure and law forward and transport the
likelihood by \(S^{-1}\). & Bayes-factor comparisons are preserved on
corresponding observations, with the density-version and almost-everywhere
qualifications of \cref{thm:p15-model-density-invariance}. A newly flat
parameter prior is not the transported old prior. \\
Subset mass assignment and additive allocations &
\(m\mapsto\mathcal C_m\); distribute each positive subset mass by a
probability mass supported inside that subset. & Event infima and suprema
equal \(\operatorname{Bel}_m\) and \(\operatorname{Pl}_m\), with attained
bounds, \cref{thm:p14-evidence-finite-measures}. General lower envelopes
need not have this mass form, \cref{thm:p14-evidence-measure-separation}. \\
Finite prior list and its convex hull &
Replace \(\mathcal C=\{p^1,\ldots,p^m\}\) by its probability-vector
convex hull. & Every act's minimum expected utility is unchanged,
\cref{thm:p17-finite-maxmin-properties}. The two prior sets need not be
the same set, and a probability distribution over their members is not
provided by taking the hull. \\
\bottomrule
\end{longtable}

These rows preserve specified numerical functions on their declared
objects. They are not embeddings of every uncertainty logic into one
another. In particular, max-based possibility, subset-mass belief and an
additive event probability have distinct operations and interpretations;
the complete finite comparison is \cref{tab:p14-evidence-eight-events}.
Fuzzy truth values give still another semantic role for numbers in
\([0,1]\), fixed by the named algebra and designated-value conditions.
Changing that role requires a justified bridge.

\section{Separating examples and limits of transfer}
\label{sec:p14-translation-separating-limits}

\subsection{A first-order property outside basic modal expression}

Let the first pointed frame have \(W=\{a,b\}\) and only edge \(aRb\),
and the second have \(W'=\{a',b',c'\}\) and only edges
\(a'R'b'\), \(a'R'c'\). Make every propositional atom false at every
world. The relation
\[
 Z=\{(a,a'),(b,b'),(b,c')\}
\]
is a bisimulation. At the root pair each successor has a match in both
directions; the leaf pairs have no successors and agree on all atoms.
The points therefore agree on all basic modal formulas by
\cref{thm:p08-modal-bisimulation-invariance}.

The first-order formula
\[
 \exists y\exists z\,(Rxy\land Rxz\land y\ne z)
\]
is false at \(a\) and true at \(a'\). It counts two distinct
successors, which that bisimulation does not preserve. Thus it cannot be
a basic modal standard translation. This concrete finite witness uses the
invariance direction; it needs no finite-model converse characterization.
The relation \(Z\) is not required to be a function or an injection.

\subsection{The same finite graph under full and selected relation ranges}

\begin{hbdefinition}{def:p14-selected-reachability-model}{A two-element relation-range comparison}
Use the relational second-order signature with object constants \(a,b\),
binary predicate constant \(E\), and logical equality. Let \(M=\{0,1\}\),
interpret \(a=0,b=1,E=\varnothing\), and define the selected relation
domains by
\[
 \mathcal R_1=\{M^1\},\qquad
 \mathcal R_2=\{\varnothing,M^2\},\qquad
 \mathcal R_n=\{M^n\}\quad(n\geq3).
\]
Here a unary relation is a set of one-entry tuples; membership of \((x)\)
corresponds to membership of \(x\) in the associated subset of \(M\).
These data give an arbitrary relational Henkin structure \(\mathcal H\)
under \cref{def:p02-sol-henkin-structure}: every relation domain is
nonempty, and the predicate constant's empty interpretation belongs to
\(\mathcal R_2\). Let \(\mathcal F\) have the same object and predicate
interpretations but full relation domains.
\end{hbdefinition}

\begin{hbtheorem}{thm:p14-selected-reachability}{proved=annex:p14.translation.selected-reachability}{A finite domain does not force full relation semantics}
For the exact formula \(\operatorname{Reach}(a,b)\) of
\cref{def:p02-sol-reachability},
\[
 \mathcal H\entails\operatorname{Reach}(a,b),\qquad
 \mathcal F\not\entails\operatorname{Reach}(a,b).
\]
Neither graph has a finite directed path from 0 to 1. The selected structure
also fails the eligible comprehension instance picking out 0 by
\(x=a\), so it is not a matched comprehension--choice Henkin model.
This does not contradict adequacy for that matched model class.
\end{hbtheorem}

The witness isolates the quantifier range. It changes neither the finite
object domain nor the graph relation. Thus an assertion that a finite
domain makes full and arbitrary Henkin semantics interchangeable would
fail already here. The defining closure condition for the matched model
class excludes this structure for a specified reason.

\subsection{Satisfiability, categorical equality and unrestricted valuations}

Fresh-symbol clausification offers another kind of bridge. In domain
\(\{0,1\}\) with \(P=\{0\}\), \(\exists xP(x)\) is true.
An expansion with fresh constant \(c=0\) satisfies \(P(c)\); an expansion
with \(c=1\) does not. \Cref{thm:p12-small-skolem} requires the existence
of a suitable expansion, not truth in every expansion or same-signature
equivalence. Reusing a preassigned constant would change the model relation.

In a CCC interpretation assigning a singleton to a base type, the two
context projections denoting variables \(x,y\) of that type coincide.
Assigning the two-element set in \(\mathsf{Set}\) makes them distinct:
at the pair \((0,1)\), their values are 0 and 1. Hence equality in one
interpretation is insufficient for typed \(\beta\eta\) equality; the
criterion in \cref{thm:p10-stlc-ccc-completeness} requires every
interpretation. The context products and terminal factors are those of
\cref{def:p10-stlc-ccc-interpretation}, not extra term constructors.

Finally, the Boolean-operator embedding preserves evaluations whose
assigned elements lie in its clopen image. It does not automatically
preserve every equation under all valuations in the full powerset
completion. General-frame completeness and ordinary-frame completeness
therefore require the different hypotheses in
\cref{thm:p10-modal-general-frame-strong,
thm:p10-modal-canonicity-completeness}. The existence of Kripke-incomplete
normal logics, cited in \cref{thm:p08-kripke-incomplete-logics}, is a
separate reason to retain that distinction.

\begin{longtable}{@{}p{.27\textwidth}p{.31\textwidth}p{.32\textwidth}@{}}
\caption{Satisfiability, categorical equality and unrestricted valuations}\\
\toprule
Transfer being attempted & Decisive limit & Correct interface to retain \\
\midrule
\endfirsthead
\caption[]{Satisfiability, categorical equality and unrestricted valuations (continued)}\\
\toprule
Transfer being attempted & Decisive limit & Correct interface to retain \\
\midrule
\endhead
\midrule\multicolumn{3}{r}{\emph{continued on next page}}\\
\endfoot
\endlastfoot
Every FO relation property is modal & Distinct-successor counting changes
across the displayed bisimulation. & Standard-translation image and the
bisimulation-invariance condition. \\
A propositional Glivenko row is a quantified translation theorem & Its
source statement restricts both calculi to propositional rules. & Use the
explicit first-order recursive map and its stated context transformation. \\
Full categoricity combines with Henkin completeness & The two claims
quantify over different relation/function ranges. & Full-model and
matched-general-model statements separately. \\
Finite domains remove relation-range conditions & The two-element
reachability witness above. & Keep each \(\mathcal R_n\), its admitted
predicate constants and schema closure. \\
An incidence translation proves native relational-rule completeness & The
calculus is defined by its reduction; admissible projections alone give no
converse generation theorem. & \cref{def:p02-sol-native-boundary}. \\
Derivability equivalence preserves proof normalization & The nonvacuous
translation inserts the obstructed detour. &
\cref{thm:p07-nj-vacuous-consequence-bridge,
thm:p07-nj-nonvacuous-obstruction}. \\
An object or subalgebra embedding supplies every target valuation & Image
valuations form a restricted class. & Clopen/admissible representation;
additional canonicity for unrestricted frames. \\
Unrestricted interpolation or preservation transfers to finite structures &
The relevant compactness and finite-class hypotheses change. &
\cref{thm:p09-finite-craig-beth,thm:p03-los-tarski}. \\
\bottomrule
\end{longtable}

Every row is a comparison under declared assumptions. Formula-image size,
effective construction, proof-length growth and equivalence of algorithms
require further statements if they matter to the intended use. The source
and proof pointers give the exact scope from which a transfer can be made;
an embedding name alone supplies none of those additional conclusions.
